%=====================================================================
\section{Thermodynamic cycles and the energy budget}
\label{sec:cycles}

%---------------------------------------------------------------------
\subsection{State points}
\label{sec:cycles:states}

The plant contains two cycles. During charging a two-stage high-temperature heat
pump with two regenerators lifts heat from a source at $T_{\mathrm{source}}$ to
the PCM melting temperature. During discharging a two-stage organic Rankine
cycle rejects to a sink at $T_{\mathrm{sink}}$. Both use cyclopentane; the
secondary fluid circulating through the boreholes is pressurised water.

Every state point is fixed by a specified approach temperature rather than by a
component model. Writing $\Tm$ for the PCM melting point,
\begin{align}
  T_{1e} &= T_{\mathrm{sink}} + \Delta T_{\mathrm{sink,glide}}
            + \Delta T_{E,\mathrm{sink}}, &
  T_{3d} &= \Tm^{\mathrm{bot}} - \Delta T_{M,1D}, \\
  T_{2d} &= T_{3d} + \Delta T_{3D,2D}\ \text{(provisional)}, &
  T_{3e} &= \min\!\left(T_{2d}-\Delta T_{2D,3E},\,
                        T_{3e}^{\mathrm{pinch}}\right), \\
  T_{4c} &= \Tm + \Delta T_{4C,M}, &
  T_{4a} &= T_{\mathrm{source}} - \Delta T_{3A,4A}, \\
  T_{13h} &= T_{4a} - \Delta T_{\mathrm{pinch,HPE}}, &
  T_{3c} &= T_{4c}, \\
  T_{2c} &= T_{3c} - \Delta T_{3C,2C}, &
  T_{2h} &= T_{3c} + \Delta T_{2H,3C},
\end{align}
with
\begin{equation}
  \Tm^{\mathrm{bot}} \;=\; \Tm - \Delta T_{\mathrm{cascade}}\,
  \frac{N_{\mathrm{lay}}-1}{N_{\mathrm{lay}}}
  \label{eq:Tmbot}
\end{equation}
the melting temperature of the coldest cascade layer, \SI{101.111}{\celsius} for
the default case. Function: \code{T\_m\_bottom}, \code{cycle\_state\_points}.

Three of these approaches are \emph{derived} rather than free, and saying so is
the point of the arrangement above. $T_{13h}$ follows from the source outlet
(\S\ref{sec:defects:hpe}); $T_{3e}$ is the smaller of two constraints
(\S\ref{sec:cycles:pinch}); and $T_{1e}$ is referenced to the sink
\emph{outlet}, which coincides with its inlet only because
$\Delta T_{\mathrm{sink,glide}}=0$ models the sink as a reservoir.

The charging glide $\Delta T_{3C,2C}$ is the principal design variable of the
study, and since v0.6 its only job is to set the two mass flow rates. Until
v0.9b the discharging glide and the cascade grading were not separate fields at
all: both were $\Delta T_{3C,2C}$. They are now
$\Delta T_{3D,2D}$ and $\Delta T_{\mathrm{cascade}}$, each defaulting to it, and
\S\ref{sec:cycles:glides} sets out which of the resulting equalities are forced
by energy conservation and which were merely assumed.

\subsection{Which temperatures may be prescribed}
\label{sec:cycles:closure}

The arrangement above changed in v0.6, and the reason is worth stating because
it is easy to get wrong in either direction.

\emph{Only the two exchanger inlets are boundary conditions.} $T_{4c}$ and
$T_{3d}$ are what the marches of \S\ref{sec:front:enthalpy} are given. Both
outlets --- $T_{2c}$ back to the heat-pump condenser, $T_{2d}$ on to the ORC
evaporator --- are computed by the model, and cannot also be specified.

Until v0.5a the discharge was closed the other way round: the \emph{outlet} was
pinned at $T_{2d}=\Tm-\Delta T_{m,2D}$ with $\Delta T_{m,2D}=0$, and the inlet
obtained by subtracting the glide. Two objections, the second fatal.

First, it is contradicted by the solution. Under Formulation~C the PCM leaves
$\Tm$, and at cyclic steady state the store finishes each charge superheated, so
at the start of the following discharge the water leaves the field at
\SI{157.93}{\celsius} --- nearly \SI{8}{\kelvin} above the melting temperature of
the top layer, the very quantity the closure had pinned it to. The same applies
on charge: $T_{2c}$ was prescribed at \SI{105.00}{\celsius} and the model
returns a flow-mixed mean of \SI{104.38}{\celsius}, ranging over more than
\SI{20}{\kelvin} within the window.

Second, the implied inlet approach was set by nobody. With $T_{2d}=\Tm$ and a
glide of $\Delta T_{3C,2C}$ across $N_{\mathrm{lay}}$ equal layers, the
discharge entered the coldest layer at
\begin{equation}
  \Delta T_{\mathrm{implied}}
  = \frac{\Delta T_{3C,2C}}{N_{\mathrm{lay}}}
  = \SI{6.111}{\kelvin},
  \label{eq:implied-approach}
\end{equation}
exactly one layer width --- and therefore changing silently with the number of
layers. Against the \SI{10}{\kelvin} the charge receives, the discharge was
being given \SI{61}{\percent} of the driving difference at every layer, and in
the linear-glide idealisation the fluid arrives at each layer's exit face
\emph{at} that layer's melting temperature, pinching to zero $N_{\mathrm{lay}}$
times over. The charge never comes closer than \SI{3.889}{\kelvin}.

The replacement is symmetric and says what it means:
\begin{equation}
  \underbrace{T_{4c}=\Tm+\Delta T_{4C,M}}_{\text{superheat, drives melting}},
  \qquad
  \underbrace{T_{3d}=\Tm^{\mathrm{bot}}-\Delta T_{M,1D}}
    _{\text{subcooling, drives freezing}},
  \label{eq:inlets}
\end{equation}
with $\Delta T_{M,1D}=\Delta T_{4C,M}=\SI{10}{\kelvin}$ by default. State 1d is
the borehole inlet on discharge in the plant diagram. Setting
$\Delta T_{M,1D}=\Delta T_{3C,2C}/N_{\mathrm{lay}}$ recovers the retired closure
to the last bit, which is how the change was verified to be a reparameterisation
rather than a new model.

\emph{What it costs.} $T_{3e}$, and hence $\eta_{\mathrm{ORC}}$, follows
$T_{2d}$. Demoting $T_{2d}$ to a result therefore couples the plant to the
store. It is handled by one correction pass rather than an outer loop, for a
reason that is checked rather than assumed --- see \S\ref{sec:cyclic:closure}.

%---------------------------------------------------------------------
\subsection{What actually sets the evaporating temperature}
\label{sec:cycles:pinch}

$T_{3e}$ does \emph{not} follow from $T_{2d}$ by a hot-end approach, although
that is how it was written until v0.9 and it is the natural thing to write. The
binding constraint lives in the middle of the evaporator, and the reason is
structural rather than numerical.

The borehole water is a \emph{sensible} stream: it glides the full
$\Delta T_{3C,2C}$, by construction, because that glide is what sizes the flow.
The ORC working fluid is a \emph{latent} stream over most of its duty: at the
design point it absorbs \SI{76}{\percent} of its heat boiling at the single
temperature $T_{3e}$. A sloping line and a horizontal one converge, and where
they meet is set by the whole exchanger, not by either end:
\begin{equation}
  \Delta T_{\mathrm{pinch}} \;=\;
  \min_{0 \le Q \le Q_{\mathrm{tot}}}
  \left[\,T_{\mathrm{water}}(Q) - T_{\mathrm{ORC}}(Q)\,\right]
  \;\ge\; \Delta T_{\mathrm{pinch,ORC}} .
  \label{eq:pinch}
\end{equation}
Both composites are built on real enthalpy --- the water because $c_{p}$ is not
constant over \SI{55}{\kelvin}, the working fluid because the latent plateau
must appear as a vertical segment and not be smeared into a slope. The cold
composite carries two streams, the main flow from state $10e$ to $3e$ at
$p_{\mathrm{evap}}$ and the fraction $y$ reheated from $5e$ to $6e$ at
$p_{\mathrm{int}}$.

Equation~\eqref{eq:pinch} is what determines $T_{3e}$: the hot-end rule is kept
when it already satisfies it, and $T_{3e}$ is otherwise lowered until it does.
At the design point the hot-end rule gives \SI{134.1}{\celsius} and
Eq.~\eqref{eq:pinch} permits only \SI{95.4}{\celsius}, which is not a small
correction and is worked through in \S\ref{sec:defects:orcpinch}.

\emph{The general lesson, which matters more than the number.} Any storage
concept that delivers heat as a \emph{glide} is badly matched to a power cycle
that absorbs it at \emph{one temperature}, and the mismatch tightens as the
glide widens. The glide cannot be treated as a free lever for storage density.
Functions: \code{orc\_pinch}, \code{feasible\_rankine}.

\medskip
\noindent\textbf{Two constraints, not one.} $T_{3e}$ is the smaller of a
minimum approach at the hot end and the largest pinch-feasible value,
\begin{equation}
  T_{3e} \;=\; \min\!\left(T_{2d}-\Delta T_{2D,3E},\;\;
                            T_{3e}^{\mathrm{pinch}}\right),
  \label{eq:T3e-min}
\end{equation}
and neither subsumes the other. At $\Delta T_{2D,3E}=\SI{12}{\kelvin}$ the
pinch binds across the whole useful range of glides, \SIrange{15}{70}{\kelvin},
so the hot-end rule is inert in practice while remaining a real constraint if
raised. The active one is reported.

%---------------------------------------------------------------------
\subsection{The three glides, and which of them are actually independent}
\label{sec:cycles:glides}

Three temperature spans in this plant are close to one another in magnitude and
were, until v0.9b, all set by the single field $\Delta T_{3C,2C}$. They are not
equally free, and the distinction matters for any parametric study.

\begin{center}
\begin{tabular}{@{}l l l@{}}
\toprule
span & value & status \\
\midrule
HTHP condenser water   & \SI{55.0}{\kelvin} & \emph{specified}: $\Delta T_{3C,2C}$ \\
borehole charging      & \SI{55.0}{\kelvin} & \textbf{forced} equal to the above \\
ORC evaporator water   & \SI{55.0}{\kelvin} & \emph{specified}: $\Delta T_{3D,2D}$ \\
borehole discharging   & \SI{55.0}{\kelvin} & \textbf{forced} equal to the above \\
PCM cascade span       & \SI{48.889}{\kelvin} & \emph{specified}: $\Delta T_{\mathrm{cascade}}$ \\
\bottomrule
\end{tabular}
\end{center}

\noindent
The two \textbf{forced} equalities are not modelling assumptions and cannot be
relaxed. The state-point construction sets $T_{3c}=T_{4c}$: the water leaves the
heat-pump condenser and enters the borehole with no state change between, so it
is the same stream at the same flow. Its rise in the condenser and its fall in
the borehole are the same number by energy conservation. The discharge side is
identical in structure.

What \emph{was} assumed is only that the charging and discharging glides are
equal --- both mass flows were divided by $\Delta T_{3C,2C}$ --- and that the
cascade is graded to match them.

\medskip
\noindent\textbf{The cascade span, and why its bound runs the way it does.}
With a span one layer short of the glide, the approach between the water and
the layer it is working on is \emph{exactly} $\Delta T_{4C,M}$ at the leading
face of every layer, decaying to $\Delta T_{4C,M}-\Delta T_{\mathrm{cascade}}/N_{\mathrm{lay}}$
at each trailing face --- \SI{10.000}{\kelvin} and \SI{3.889}{\kelvin} at the
design point. That sawtooth is the cascade, and matching the span to the glide
is the \emph{unique} choice that makes it uniform.

The span is therefore bounded \textbf{below}, not above:
\begin{equation}
  \mathrm{span} \;>\; \max\!\left(
    \Delta T_{3C,2C}-\Delta T_{4C,M},\;\;
    \Delta T_{3D,2D}-\Delta T_{M,1D}\right) = \SI{45.0}{\kelvin},
  \label{eq:span-bound}
\end{equation}
against a design value of \SI{48.889}{\kelvin}: a margin of only
\SI{3.889}{\kelvin}. Below the bound the driving difference \emph{inverts} at
the far end of the well, and \code{validate\_case} raises. The narrow direction
is the dangerous one, and it is also the tempting one, because narrowing the
cascade raises $T_{m,\mathrm{bottom}}$ and lets the ORC boil hotter:

\begin{center}
\begin{tabular}{@{}r r r r r@{}}
\toprule
span [\si{\kelvin}] & $T_{m,\mathrm{bot}}$ [\si{\celsius}] & deviation
 & $\varepsilon$ end dc & $\eta_{\mathrm{RTE}}$ \\
\midrule
26.7 & 123.3 & \SI{-32.56}{\percent} & 0.196 & \textbf{0.3508} \\
35.6 & 114.4 & \SI{-17.19}{\percent} & 0.108 & 0.3335 \\
43.5 & 106.5 & \SI{-4.00}{\percent}  & 0.041 & 0.3149 \\
\textbf{48.9} & \textbf{101.1} & \SI{+4.53}{\percent} & \textbf{0.012} & 0.3003 \\
\bottomrule
\end{tabular}
\end{center}

\noindent
Efficiency and utilisation pull in opposite directions, so there is an interior
optimum and it has not been searched.

\medskip
\noindent\textbf{The discharging glide.} $\Delta T_{3D,2D}$ shows the same
structure. A wider discharge glide raises $T_{2d}$, lets the ORC boil hotter
and lifts $\eta_{\mathrm{RTE}}$ --- until the store cannot keep up:

\begin{center}
\begin{tabular}{@{}r r r r@{}}
\toprule
$\Delta T_{3D,2D}$ [\si{\kelvin}] & deviation & $\varepsilon$ end dc
 & $\eta_{\mathrm{RTE}}$ \\
\midrule
40 & \SI{+6.41}{\percent} & 0.000 & 0.2764 \\
55 & \SI{+4.53}{\percent} & 0.012 & 0.3003 \\
65 & \SI{-8.51}{\percent} & 0.227 & \textbf{0.3041} \\
\bottomrule
\end{tabular}
\end{center}

\noindent
\emph{Why the three were tied together in the first place.} The cascade is
shared hardware: graded once, and both half-cycles must live with it. Matching
both water glides to the span is the only way to hold a uniform approach in
both directions. That is a defensible design choice, and it is now expressible
as a choice rather than embedded as an identity.

Given the state points, the cycle efficiencies
\begin{equation}
  \eta_{\mathrm{ORC}} \;=\; \frac{w_{\mathrm{net}}}{q_{\mathrm{in}}},
  \qquad
  \mathrm{COP} \;=\; \frac{q_{\mathrm{out}}}{w_{\mathrm{comp}}}
  \label{eq:cycle-eff}
\end{equation}
follow from enthalpies evaluated by CoolProp at those points.

\emph{Source:} \code{\_legacy\_physics.double\_stage\_rankine},
\code{two\_stage\_htheatpump\_2regs}; assembled in \code{system.\_cycles}.

\medskip
\noindent\textbf{Caveat, stated here because it bounds every efficiency in this
document.} The expansion and compression steps are evaluated between the
prescribed state points without isentropic efficiencies. $\eta_{\mathrm{ORC}}$
and $\mathrm{COP}$ are therefore idealised, and the round-trip efficiencies of
\S\ref{sec:hydraulics:kpi} inherit that idealisation. Only the electrical and
mechanical efficiencies $\eta_{T}$, $\eta_{G}$, $\eta_{C}$, $\eta_{H}$ of
Eq.~\eqref{eq:budget} appear. This is a property of the model as built, not a
simplification introduced by this document.

%---------------------------------------------------------------------
\subsection{Energy budget}
\label{sec:cycles:budget}

The chain runs backwards from the specified net electrical output
$\dot W_{\mathrm{el,out}}$ of the power block:
\begin{align}
  \dot W_{T} &= \frac{\dot W_{\mathrm{el,out}}}{\eta_{T}\,\eta_{G}}, &
  \dot Q_{\mathrm{in,ORC}} &= \frac{\dot W_{T}}{\eta_{\mathrm{ORC}}},
  \nonumber\\[2pt]
  \Delta E_{\mathrm{in,ORC}} &= \dot Q_{\mathrm{in,ORC}}\, t_{\mathrm{dc}}, &
  \Delta E_{\mathrm{out,HP}} &= \Delta E_{\mathrm{in,ORC}}\,(1+\lambda),
  \nonumber\\[2pt]
  \dot Q_{\mathrm{out,HP}} &= \frac{\Delta E_{\mathrm{out,HP}}}{t_{\mathrm{ch}}}, &
  \dot W_{C,\mathrm{HP}} &= \frac{\dot Q_{\mathrm{out,HP}}}{\mathrm{COP}},
  \nonumber\\[2pt]
  \dot W_{\mathrm{el,in}} &= \frac{\dot W_{C,\mathrm{HP}}}{\eta_{C}\,\eta_{H}}. &&
  \label{eq:budget}
\end{align}

The charging mass flow of secondary fluid follows from the duty and the glide,
\begin{equation}
  \mdw^{\mathrm{ch}}
  \;=\;
  \frac{\dot Q_{\mathrm{out,HP}}}{c_{p,w}\,\Delta T_{3C,2C}},
  \qquad
  c_{p,w} = c_p\!\left(\tfrac{1}{2}(T_{3c}+T_{2c}),\,P\right),
  \label{eq:mdot-charge}
\end{equation}
and is divided equally among all tubes of all wells,
\begin{equation}
  \md_{1}^{\mathrm{ch}} \;=\; \frac{\mdw^{\mathrm{ch}}}{\Nwells\, n_{t}} .
  \label{eq:mdot-per-tube}
\end{equation}

Note for \S\ref{sec:hydraulics:sizing}: Eq.~\eqref{eq:mdot-charge} \emph{pins}
the charging flow to the specified glide. It is therefore not free, and this is
what decides which side of the cycle can size the well field.

%---------------------------------------------------------------------
\subsection{The storage-loss surplus \texorpdfstring{$\lambda$}{lambda}}
\label{sec:cycles:lambda}

The factor $(1+\lambda)$ in Eq.~\eqref{eq:budget} is the one place where the
budget anticipates a loss it does not compute. It asserts that charging must
deposit $\lambda$ more energy than discharging will recover, and it is set to
$\lambda = 0.05$ by fiat.

In v0.1 this could not be checked, because the discharge model began from a bare
tube in fully solid PCM irrespective of what charging had produced
(\S\ref{sec:front:legacy}). With the melt state carried across the cycle
(\S\ref{sec:front:marched}) the recovered fraction becomes an output,
\begin{equation}
  \eta_{\mathrm{storage}}
  \;=\;
  \frac{\Delta E_{\mathrm{discharged}}}{\Delta E_{\mathrm{stored}}},
  \label{eq:eta-storage}
\end{equation}
and can be compared with the $1/(1+\lambda)=0.952$ that Eq.~\eqref{eq:budget}
presumes. At the design point with the legacy wall conductivity
$\eta_{\mathrm{storage}}=0.953$, and the assumption survives. With the correct
steel conductivity it falls to $0.9364$, and it does not: the field banks about
\SI{7}{\percent} more than the ORC withdraws, against the \SI{5}{\percent}
assumed. Making $\lambda$ an output rather than an input remains open work.
